% GENERATED by scripts/generate-mindmaps.mjs — do not hand-edit.
\begin{tikzpicture}[
  mindmap,
  every node/.style={
    concept,
    execute at begin node=\hskip0pt,
    align=center,
    font=\sffamily
  },
  mm root/.style={
    concept color=graphite,
    text=ghostwhite,
    font=\sffamily\bfseries,
    minimum size=3.35cm,
    text width=2.75cm
  },
  mm section/.style={
    font=\sffamily\bfseries\scriptsize,
    minimum size=2.42cm,
    text width=1.92cm
  },
  mm topic/.style={
    font=\sffamily\tiny,
    minimum size=1.92cm,
    text width=1.46cm
  }
]
  \node[mm root] (root) at (0,0) {Chapter 9\\Deterministic Omega-Automata and Star-Free Languages};
  \node[mm section, concept color=brightsky] (s1) at (68.82:5.1cm) {The Quest for Determinism in Infinite Verification};
  \node[mm topic, concept color=brightsky!44!white] (t1_1) at (79.41:9.8cm) {An Intuitive Example: Non-Deterministic Guessing vs. Real-Time Control};
  \node[mm topic, concept color=brightsky!44!white] (t1_2) at (58.24:9.8cm) {Why the Subset Construction Fails on Büchi Automata};
  \node[mm section, concept color=vvteal] (s2) at (15.88:5.1cm) {Deterministic Büchi Automata (DBA) and the Limit of Vectorial Closure};
  \node[mm topic, concept color=vvteal!44!white] (t2_1) at (37.06:9.8cm) {Vectorial Closure};
  \node[mm topic, concept color=vvteal!44!white] (t2_2) at (15.88:9.8cm) {The DBA Characterization Theorem};
  \node[mm topic, concept color=vvteal!44!white] (t2_3) at (-5.29:9.8cm) {The Limit of DBAs: The Finitely Many a's Language};
  \node[mm section, concept color=vvviolet] (s3) at (-68.82:5.1cm) {Deterministic Muller Automata (DMA) and McNaughton's Theorem};
  \node[mm topic, concept color=vvviolet!44!white] (t3_1) at (-26.47:9.8cm) {Muller Acceptance Condition};
  \node[mm topic, concept color=vvviolet!44!white] (t3_2) at (-47.65:9.8cm) {Closure Properties of DMAs};
  \node[mm topic, concept color=vvviolet!44!white] (t3_3) at (-68.82:9.8cm) {The DMA Characterization Lemma};
  \node[mm topic, concept color=vvviolet!44!white] (t3_4) at (-90:9.8cm) {McNaughton's Theorem and the Closure Bottleneck};
  \node[mm topic, concept color=vvviolet!44!white] (t3_5) at (-111.18:9.8cm) {Corollary: Weak S1S (WS1S) = S1S};
  \node[mm section, concept color=vvorange] (s4) at (-153.53:5.1cm) {Deterministic Rabin Automata, Parity, and Safra's Construction};
  \node[mm topic, concept color=vvorange!44!white] (t4_1) at (-132.35:9.8cm) {Rabin Acceptance Condition};
  \node[mm topic, concept color=vvorange!44!white] (t4_2) at (-153.53:9.8cm) {Parity Acceptance and the Equivalence DMA \ensuremath{\equiv} DRA \ensuremath{\equiv} Parity};
  \node[mm topic, concept color=vvorange!44!white] (t4_3) at (-174.71:9.8cm) {Safra's Determinization (Roadmap)};
  \node[mm section, concept color=tomato] (s5) at (-227.65:5.1cm) {The First-Order Fragment of S1S and Star-Free Languages};
  \node[mm topic, concept color=tomato!44!white] (t5_1) at (-195.88:9.8cm) {The Signature FO-S1S[<]};
  \node[mm topic, concept color=tomato!44!white] (t5_2) at (-217.06:9.8cm) {Star-Free Languages};
  \node[mm topic, concept color=tomato!44!white] (t5_3) at (-238.24:9.8cm) {Translation: Star-Free \ensuremath{\Rightarrow} FO-S1S[<]};
  \node[mm topic, concept color=tomato!44!white] (t5_4) at (-259.41:9.8cm) {Translation: FO-S1S[<] \ensuremath{\Rightarrow} Star-Free (McNaughton-Papert)};

  \begin{scope}[on background layer, line cap=round]
    \draw[brightsky, line width=11pt] (root) -- (s1);
    \draw[brightsky!60!white, line width=5.5pt] (s1) -- (t1_1);
    \draw[brightsky!60!white, line width=5.5pt] (s1) -- (t1_2);
    \draw[vvteal, line width=11pt] (root) -- (s2);
    \draw[vvteal!60!white, line width=5.5pt] (s2) -- (t2_1);
    \draw[vvteal!60!white, line width=5.5pt] (s2) -- (t2_2);
    \draw[vvteal!60!white, line width=5.5pt] (s2) -- (t2_3);
    \draw[vvviolet, line width=11pt] (root) -- (s3);
    \draw[vvviolet!60!white, line width=5.5pt] (s3) -- (t3_1);
    \draw[vvviolet!60!white, line width=5.5pt] (s3) -- (t3_2);
    \draw[vvviolet!60!white, line width=5.5pt] (s3) -- (t3_3);
    \draw[vvviolet!60!white, line width=5.5pt] (s3) -- (t3_4);
    \draw[vvviolet!60!white, line width=5.5pt] (s3) -- (t3_5);
    \draw[vvorange, line width=11pt] (root) -- (s4);
    \draw[vvorange!60!white, line width=5.5pt] (s4) -- (t4_1);
    \draw[vvorange!60!white, line width=5.5pt] (s4) -- (t4_2);
    \draw[vvorange!60!white, line width=5.5pt] (s4) -- (t4_3);
    \draw[tomato, line width=11pt] (root) -- (s5);
    \draw[tomato!60!white, line width=5.5pt] (s5) -- (t5_1);
    \draw[tomato!60!white, line width=5.5pt] (s5) -- (t5_2);
    \draw[tomato!60!white, line width=5.5pt] (s5) -- (t5_3);
    \draw[tomato!60!white, line width=5.5pt] (s5) -- (t5_4);
  \end{scope}
\end{tikzpicture}
